\usepackage{iftex}
\usepackage{needspace}
\ifluatex
  % Fonts with the coverage this paper needs, loaded only under lualatex:
  % a monospace with IPA for code, DejaVu Sans for IPA in tables, a Korean
  % face for the Hangul sound tokens, and a Chinese face (with CJK
  % Extension A) for the cluster tokens.
  \usepackage{fontspec}
  \setmonofont{DejaVu Sans Mono}[Scale=MatchLowercase]
  \newfontfamily\ipafont{Noto Sans}[Scale=MatchLowercase]
  \newfontfamily\tokenfont{Apple SD Gothic Neo}[Scale=MatchLowercase]
  \newfontfamily\clusterfont{Hiragino Sans GB}[Scale=MatchLowercase]
\else
  \newcommand{\ipafont}{}
  \newcommand{\tokenfont}{}
  \newcommand{\clusterfont}{}
\fi
\usepackage[english]{babel}
\usepackage{graphicx}
\usepackage{amsmath}
\usepackage{amssymb}
\usepackage{amsfonts}
\usepackage[thmmarks, amsmath]{ntheorem}
\usepackage{setspace}
\usepackage{csquotes}
\usepackage{enumitem}
\usepackage{parskip}
\usepackage{microtype}
\usepackage[htt]{hyphenat}
\usepackage{booktabs}
\usepackage{array}
\usepackage{tabularx}
\usepackage{longtable}
\usepackage{fancyvrb}
\usepackage{xcolor}
\usepackage[margin=1in]{geometry}
\usepackage{fancyhdr}
\usepackage{titlesec}
\usepackage[
    pdfauthor={Lance Pollard},
    pdftitle={Talk: A Phonetic Encoding for Speech and Language Models},
    pdfsubject={A phonetic encoding with an ASCII form and a one-token-per-sound machine form, for speech and language models},
    pdfkeywords={phonetics, IPA, phonetic transcription, tokenization, NLP, speech, text-to-speech, grapheme-to-phoneme, romanization, X-SAMPA, subword, Hangul},
    pdfcreator={Lance Pollard},
    pdfproducer={LuaLaTeX},
    colorlinks=true,
    linkcolor=black,
    urlcolor=black,
    citecolor=black
]{hyperref}

\hypersetup{
    pdfinfo={
        ContactEmail={lancejpollard@gmail.com},
        ContactURL={https://clue.surf},
        Copyright={Copyright (C) 2026 ClueSurf}
    }
}

\usepackage{etoolbox}
\graphicspath{{figure/}}
\apptocmd{\thebibliography}{\interlinepenalty=10000}{}{}

\widowpenalty=10000
\clubpenalty=10000
\setlength{\parindent}{0pt}
\setlength{\parskip}{1em}
\tolerance=1500
\emergencystretch=3em
\setlength{\parfillskip}{0pt plus 1fil}
\sloppy

\pagestyle{fancy}
\fancyhf{}
\renewcommand{\headrulewidth}{0pt}
\fancyfoot[C]{-- \thepage\ --}
\setlength{\footskip}{40pt}
\fancypagestyle{plain}{
  \fancyhf{}
  \fancyfoot[C]{-- \thepage\ --}
  \renewcommand{\headrulewidth}{0pt}
}

\BeforeBeginEnvironment{equation}{\vspace{4pt}}
\AfterEndEnvironment{equation}{\vspace{4pt}}
\setlength{\itemsep}{0pt}
\setlist[itemize]{before=\vspace{0pt},after=\vspace{-2pt}}

\titleformat{\section}{\Large\bfseries}{\thesection.}{0.5em}{}
\titleformat{\subsection}{\normalsize\bfseries}{\thesubsection.}{0.5em}{}
\titleformat{\subsubsection}{\small\itshape}{\thesubsubsection.}{0.5em}{}

\onehalfspacing

\theoremheaderfont{\bfseries}
\theorembodyfont{\normalfont}
\theoremindent0pt
\newtheorem{definition}{Definition}
\newtheorem{principle}{Principle}

\makeatletter
\def\@begintheorem#1#2#3{%
  \par\addvspace{1em}%
  \noindent\textbf{#1 #2.}
  \ifx\@empty#3\relax
    \quad
  \else
    \ \textbf{#3:}
  \fi
  \hspace{0.5em}
}
\def\@endtheorem{\par}
\makeatother

% Inline literal talk spellings and code, robust to the special
% characters talk uses (~ ! @ ^ _ & $ *).
\newcommand{\code}[1]{\texttt{#1}}

% Code blocks use fancyvrb Verbatim with commandchars, so comments can be
% wrapped in a lighter gray. IPA glyphs render through the font, unlike
% listings, which mishandles multi-byte input. Use \cmt{...} for a comment.
\definecolor{codecomment}{gray}{0.55}
\newcommand{\cmt}[1]{\textcolor{codecomment}{#1}}

\title{
  {\huge \textbf{Talk}} \\[3mm]
  \large A phonetic encoding with one token per effective sound, \\ for speech and language models
}
\author{Lance Pollard \\ \small lancejpollard@gmail.com}
\date{July 2026}
